\section{Revision history and response to the independent review}\label{app:revision}
\paragraph{History.} Version 1 (October 25, 2026) reported the experiments of Appendix~\ref{app:original}. An independent
technical review (October 26, 2026) recommended major revision with thirteen findings. Version 2 answered them; the
full response, finding by finding, with the checks that confirmed or qualified each one, is \texttt{RESPONSE-TO-REVIEW.md} in
the release. The revision plan was committed before any revision run (\texttt{PLAN-REVISION.md}); its deviations R1--R\RVnDev{}
are recorded in \texttt{DEVIATIONS-PAPER.md}. Version 3 (October 26, 2026) answered a second review (eight findings;
\texttt{RESPONSE-TO-REVIEW-2.md}, plan \texttt{PLAN-REVISION-2.md}): the pulse-fidelity experiments of
Section~\ref{sec:fidelity}, the recentred fixed-symmetry control, the width contrasts paired by seed and the revised
settings tables. Version 4 (October 27, 2026) answers a third review (five findings;
\texttt{RESPONSE-TO-REVIEW-3.md}, plan \texttt{PLAN-REVISION-3.md}, committed before its runs): the per-sample stream is
compared with the literal one paired by draw, with bootstrap bounds, in a desk rerun of the update probes that keeps the
draws; the cross-cell dependence is measured by per-pair correlations that cannot cancel; the training comparison reports
both arms' labels and every rank reversal; the open training rerun of the one material cell was done; and an interval that
includes the width margin is no longer read as growth below it. Version 5 (October 27, 2026) answers a
fourth review (three findings; \texttt{RESPONSE-TO-REVIEW-4.md}, plan \texttt{PLAN-REVISION-4.md}, committed before
its runs). Version 4's five-seed rerun of the material SGD cell had joined, by seed only, two stream runs whose
initialisation and example order differed from their controls' (the random draws depend on the GPU model), and it
had put a seed-0 difference down to GPU rounding; its figures (stream gap \RYAtwoRSeqGap\,pp, change \RYAtwoRD\,pp)
are withdrawn as an unmatched, mixed-hardware diagnostic. The cell was rerun as matched triples on one GPU, and the
analysis now refuses a seed whose arms do not share both hashes. The saved pulse draws were added to the deposit inventory with a
command that rebuilds the probe tables from them, the probe table was split in two for legibility, and the abstract
states that its probe count is a point-estimate count. Version 6 (October 27, 2026) answers a fifth review
(one finding and two wording corrections; \texttt{RESPONSE-TO-REVIEW-5.md}; no experiment was run and no number
changed). The HTML edition of versions 4 and 5 was incomplete: it lacked Sections~\ref{sec:energy}--\ref{sec:availability}, the
contribution and acknowledgment statements and two appendices, it printed the rows of the simulator-checks table under the
sensitivity table's caption, and nine references were unresolved. The converter had read two adjacent tables with
different column-width units as one and dropped everything between them. The conversion now rewrites paragraph-width
columns, and a new check refuses any HTML that lacks a heading or a table's own rows or that holds an unresolved reference. The
two wording corrections are in Section~\ref{sec:fidelity} and Section~\ref{sec:availability}. Version 7 (October 27, 2026)
answers a sixth review (one cosmetic finding; \texttt{RESPONSE-TO-REVIEW-6.md}; no experiment was run, and no
number or sentence of the PDF changed except this history). The HTML edition of version 6 printed the rule arguments
as text before ``Seed'' in the header of Table~\ref{tab:rev4_a2m}. The conversion now removes the table-rule
commands, and the HTML check refuses any leftover rule argument or raw command. Version 8 (October 27, 2026) answers a
seventh review, which found no actionable defect (\texttt{RESPONSE-TO-REVIEW-7.md}; no experiment was run and no number
changed). It carries two changes ordered by the author: his official name, Anthony DeMarco, replaces the name printed in versions 1 to 7,
and the contribution statement credits ChatGPT's reviews. Version 9 (October 28, 2026) corrects the deviation count R1--R6 to R1--R8, found by the deposit's clean-room regeneration.
Version 10 (October 28, 2026) gives the public archive, its DOI and its GitHub mirror in Section~\ref{sec:availability}; no number changed.
Version 11 (October 28, 2026) answers an eighth review, a code audit and reproduction from the public archive
(\texttt{RESPONSE-TO-REVIEW-8.md}; no experiment was run). Figure~\ref{fig:main} (version~1, superseded) drew two negative
LM means at $+0.05$, with their intervals shifted, and its legend overlapped an axis label; both are corrected. The
energy model now counts a classification run's optimizer steps from the rows its layers trained: the run records' sample
count includes each epoch's incomplete last batch, which training skips, so the transfer reads were overstated by
0.16\,\%. No printed number changed; the largest change in a total energy is 0.0005\,\% in the current projection
and 0.0014\,\% in version~1's. Section~\ref{sec:methods} now discloses an off-by-one in the LM sampler. The
simulator's ideal reset of a partial column selection under AIHWKIT bounds, which no result used, no longer fails, and
a new check covers it.
Version 12 (October 29, 2026) answers a ninth review, a second audit and reproduction from the public archive
(\texttt{RESPONSE-TO-REVIEW-9.md}; no experiment was run). The width study's Welch interval now takes the $t$ quantile at
the Welch--Satterthwaite degrees of freedom; versions 2 to 11 rounded them down, which the plan does not state (deviation
R9). This changes one verdict, protonic-inspired c-TTv2 on CIFAR, from inconclusive to grows (Section~\ref{sec:scale}),
and moves the other Welch intervals by at most 0.18 in their units; the paired contrast is unchanged. The width section now
says that across widths a seed shares its example order, not its initialisation; the excess correlation of
Table~\ref{tab:rev3_dep} is defined as an excess in quadrature; and Table~\ref{tab:laneme} labels its version-1 values as
superseded.
Version 13 (October 31, 2026) answers a tenth review, of version 12 and version 1.2 of the archive
(\texttt{RESPONSE-TO-REVIEW-10.md}; no experiment was run and no number changed). It found the Welch correction right
and four small slips outside the results, all corrected: the archive's README said that no other printed number changed,
although all eight Welch intervals moved, and counted eight reviews; the caption of Table~\ref{tab:rev_width} dated the
rounded degrees of freedom to versions 1 to 11, not 2 to 11 (version 1 used another width test); and a self-test of the
name check left a PDF open, so its clean-up failed on Windows.
Version 14 (October 31, 2026; this paper) answers an eleventh review, of version 13 (\texttt{RESPONSE-TO-REVIEW-11.md}; no
experiment was run and no result changed). It cites Byun et al.~\cite{byun2024} and states what this study adds to it;
prints the GPU energy equation with both input ranges (the low end, \RVenGpuMeasLo\,pJ, is 500\,W at 43\,\%; version 13
named only the utilisation range); says which digital costs of the TT family are inside the energy ratio and that the ratio
is a per-operation, read-dominated figure; calls the sensitivity sheet an exploratory device sensitivity analysis, in the
subtitle too, and states that one-at-a-time ranges do not bound joint deviations; adds Figure~\ref{fig:arch}; adds a
control ladder to the limitations and computational follow-ups to the roadmap; and moves correction history from
Sections~\ref{sec:methods}, \ref{sec:fidelity}, \ref{sec:scale} and~\ref{sec:spec} to this appendix.

\paragraph{What was repaired.}
\begin{enumerate}
\item \textbf{Random streams.} The initialisation, the data order, each tile's device population and the pulse/noise draws now
      have separate generators; every record carries the hashes of its logical initialisation and data order, and every arm
      of a seed shares them (checked). The old single stream is kept as a mode that replays the original simulator bit for bit.
\item \textbf{The reset.} The original reset moved $A$ to its symmetry point and charged 62 pulses it never applied; applied
      through the device model those pulses leave \RVresetResidPct\,\% of the displacement (from $u=0.8$ the state only reaches
      \RVresetAfter, against a symmetry point of \RVresetSp; about \RVresetPulsesHundredfold{} alternating pulses would be needed
      for a hundredfold reduction). The repaired pulse reset applies and counts every pulse; the ideal reset is reported only
      as a labelled variant.
\item \textbf{The AIHWKIT presets} are mapped one to one (shifted bounds, bound-scaled slopes, slope spread), and their pulse
      responses were compared with AIHWKIT's own (Section~\ref{sec:validation}): the mapping is \RVaihwAccepted.
\item \textbf{Measured, not assumed:} the clipping of the pulse probabilities and the lost update are recorded per layer.
\end{enumerate}
The \RVnChecks{} checks of the repairs each fail on the old behaviour or a planted fault (\RVnFaults{} faults; \RVtestsAllOk{}
on the desk and on every pod).

\paragraph{How much the repairs moved the original conclusions (diagnostic D1).} Table~\ref{tab:rev_d1} runs the original and
the repaired simulator side by side at the original selections, 3 seeds, on MNIST and a 10-epoch CIFAR slice
(\RVdOneRuns{} runs).
\begin{itemize}
\item \textbf{The random streams} moved individual cells by up to \RVdOneRngMnistMaxAbs\,pp on MNIST and
      \RVdOneRngCifarMaxAbs\,pp on CIFAR (material, i.e.\ more than 1\,pp or a label change, in \RVdOneRngMnistNMaterial{} of
      \RVdOneRngMnistN{} and \RVdOneRngCifarNMaterial{} of \RVdOneRngCifarN{} cells). With three seeds this is within the noise:
      the repaired FP-SGD alone spans \RVdOneCifarFpMin--\RVdOneCifarFpMax\,\% on the CIFAR slice. The matched controls did not
      reduce the paired-gap spread here (median \RVdOneSdRev{} against \RVdOneSdOld\,pp).
\item \textbf{The reset} moved cells by up to \RVdOneResetMnistMaxAbs\,pp on MNIST and \RVdOneResetCifarMaxAbs\,pp on CIFAR. The
      ideal reset gave no hidden advantage: a physical reset (pulse, alternating pairs or none) had an equal or smaller gap in
      \RVdOneResetBetterN{} of \RVdOneResetN{} cells.
\item \textbf{The mapping} moved Li-ECRAM by at most \RVdOneMapMnistMaxAbs\,pp on MNIST and \RVdOneMapCifarMaxAbs\,pp on CIFAR.
\item \textbf{Unchanged:} plain SGD fails by far (MNIST \RVdOneSgdMnistMin--\RVdOneSgdMnistMax\,pp, CIFAR slice
      \RVdOneSgdCifarMin--\RVdOneSgdCifarMax\,pp).
\end{itemize}

\paragraph{Wording withdrawn or narrowed.} ``Near-digital'' is not used. ``Plain SGD fails on every device'' became ``every
plain-SGD arm at the FP-selected learning rate exceeded the margin''. ``The gap does not grow with width'' became the equivalence
statement of Section~\ref{sec:scale}. ``Two device properties decide'' and ``a few hundred levels suffice'' became conditional
statements about the sweeps of Section~\ref{sec:spec}. ``The memistor lacked density, speed and energy, not accuracy'' became
``an optimistic historical-device scenario attained useful accuracy in this simulator''. ``32$\times$ today with the periphery
counted'' became the conditional projection of Section~\ref{sec:energy}. ``A bench replication'' became ``a proposed, unbuilt
bench replication protocol''. ``Every matrix is analog'' became ``the learned weight matrices use the device model; attention
products and other operations are digital''.

\begin{table*}[t]\centering\scriptsize\caption{Diagnostic D1: mean gap to FP-SGD of the same simulator (pp, 3 seeds; P/F/U against 2 pp) at the original selections. Columns: orig. = the original simulator; ideal = the repaired random streams with the original ideal reset; then the reset replaced by the pulse reset (pulse), 31 alternating pulse pairs (alt.), or none; publ. = the published-style TTv2. The AIHWKIT-mapped Li-ECRAM row uses the pulse reset for TTv2 and c-TTv2. CIFAR here is a 10-epoch slice.}\label{tab:rev_d1}
\begin{tabular}{lllcccccc}\toprule Task & Device & Alg. & orig. & ideal & pulse & alt. & none & publ.\\\midrule
MNIST & Protonic-insp. & SGD & 3.09F & 4.58U &  &  &  & \\
MNIST & Protonic-insp. & TTv2 & 1.04U & 1.60U & 0.23P & 0.66P & 1.02P & 0.99P\\
MNIST & Protonic-insp. & c-TTv2 & 0.81P & 0.62U & 0.44P & 0.97U & 1.41P & \\
MNIST & Protonic-insp. & AGAD & 1.21P & 0.99P &  &  &  & \\
MNIST & Li (old map) & SGD & 4.79F & 6.47U &  &  &  & \\
MNIST & Li (old map) & TTv2 & 1.33U & 1.19U & 0.39P & 0.36P & 0.66P & 1.43P\\
MNIST & Li (old map) & c-TTv2 & 1.51U & 0.98P & 0.80P & 0.68P & 1.00P & \\
MNIST & Li (old map) & AGAD & 0.96P & 1.00P &  &  &  & \\
MNIST & Li (AIHWKIT) & SGD &  & 6.99U &  &  &  & \\
MNIST & Li (AIHWKIT) & TTv2 &  &  & 0.40P &  &  & \\
MNIST & Li (AIHWKIT) & c-TTv2 &  &  & 0.70P &  &  & \\
MNIST & Li (AIHWKIT) & AGAD &  & 0.57P &  &  &  & \\
CIFAR-10ep & Protonic-insp. & SGD & 14.24F & 10.35F &  &  &  & \\
CIFAR-10ep & Protonic-insp. & TTv2 & 5.28U & 2.53U & 2.04U & -0.83U & -1.97U & 4.17U\\
CIFAR-10ep & Protonic-insp. & c-TTv2 & 4.59U & 3.41U & 5.97U & -0.43U & 0.65U & \\
CIFAR-10ep & Protonic-insp. & AGAD & -0.70P & -2.94U &  &  &  & \\
CIFAR-10ep & Li (old map) & SGD & 17.69F & 16.91F &  &  &  & \\
CIFAR-10ep & Li (old map) & TTv2 & 4.35F & 2.51U & 2.98U & 0.57U & 0.78U & 5.41U\\
CIFAR-10ep & Li (old map) & c-TTv2 & 6.08U & 2.27U & 1.63U & 1.48U & 0.99U & \\
CIFAR-10ep & Li (old map) & AGAD & 0.39U & -1.34U &  &  &  & \\
CIFAR-10ep & Li (AIHWKIT) & SGD &  & 17.87F &  &  &  & \\
CIFAR-10ep & Li (AIHWKIT) & TTv2 &  &  & 2.37U &  &  & \\
CIFAR-10ep & Li (AIHWKIT) & c-TTv2 &  &  & 2.71U &  &  & \\
CIFAR-10ep & Li (AIHWKIT) & AGAD &  & -1.02U &  &  &  & \\
\bottomrule\end{tabular}\end{table*}


\paragraph{Version 3's update probes (superseded).} Table~\ref{tab:v3s_rev2_probe} is version 3's analysis of the update
probes (pod sitting r6). Its sampling floor for the per-sample stream added the two methods' variances as if their draws were
independent, although the per-sample stream replays the literal stream's line firings, so the floor was too large and every
excess of the per-sample stream came out as zero. Version 4 replaces it by the paired analysis of Section~\ref{sec:fidelity}
(Table~\ref{tab:rev3_probe}). The surrogate's verdicts did not depend on this floor.
The statistic of version 3, a pooled, signed normalised covariance, stays within \RYDepVthreePooledMax{} in absolute value on the same draws, because positive and negative terms cancel and unequal cell variances are pooled.

\paragraph{Random realisations of the probe and rerun states.} The probes draw from a generator of their own. On the pod (version 3) the MNIST and CIFAR probe runs reproduced their main-table records exactly; on the desk (Quadro P2200) \RYProbeRunsEqualEThree{} of \RYProbeRunsN{} do, and none shares its initialisation with its version-3 run (the MNIST run not its example order either): the simulator's generators run on the GPU, and the CUDA random kernels assign draws to threads by the GPU's multiprocessor count. The desk's probe runs are therefore fresh training realisations, and the probed states are not version 3's; the comparison within each probe is unaffected, since all three updates start from the same saved state. LM training is not bitwise reproducible even on one card type: configurations repeated exactly in the records differ by up to \RXReproLmMaxPct\,\% in bits per byte (CIFAR \RXReproCifarMaxPct\,\%).
The material SGD cell of Section~\ref{sec:fidelity} was rerun as matched triples on the desk. The pod's runs of seeds 0--2 are other random realisations of the same seeds, not a rounding variant of them: the simulator's generators run on the GPU, and the CUDA random kernels assign draws to threads by the GPU's multiprocessor count, so another GPU model draws another initialisation and example order (as predicted, the hashes differ from both the pod's and the earlier desk runs'). Per seed the two realisations' gaps differ by up to \RZRealGapDiffMax\,pp, and at the pod's three matched seeds the change was \RZPodNThreeD\,pp [\RZPodNThreeDLo, \RZPodNThreeDHi]: its gap depends strongly on the realisation.

\paragraph{The held-symmetry anchor.} With the symmetry point held at mid-range, version 2 borrowed the ordinary base (whose range is $[-1.125, +0.875]$) as this sweep's anchor and omitted its zero point (deviation R5);
version 3 ran both with each device's range re-centred (Section~\ref{sec:spec}).

% made by revision3/analyze_rev3.py from tables/rev2_probe.tex (caption prefix and label only; do not edit)
% generated by revision2/analyze_rev2.py (do not edit)
\begin{table*}[t]\centering\scriptsize\caption{Version 3 (superseded: its sampling floor for seq treated the paired draws as independent). Update fidelity (experiment A1; Li = Li-ECRAM, MO = MO-ECRAM, Prot. = protonic-inspired): one training batch applied repeatedly from the same state by the aggregate surrogate (sur.), the per-sample stream (seq) and the literal pulse-by-pulse stream (slot). Lost = unsigned requested pulse mass removed by the probability clipping (\%); $g$ = the slot stream's realised gain along the requested update; then surrogate / seq against slot: gain difference (\%), excess relative mean difference beyond the sampling floor, median per-cell variance ratio; faithful by the pre-registered lines (gain within 10 \%, excess $\le$ 0.10, variance ratio in [0.5, 2]; for seq 5 \%, 0.05, [0.8, 1.25]).}\label{tab:v3s_rev2_probe}
\begin{tabular}{lllrrrrrl}\toprule Case & Layer & Upd. & Lost & $g$ & $\Delta g$ & excess & var.\ ratio & faithful\\\midrule
CIFAR MO AGAD & 3 (128$\times$1024) & 600 & 32.1 & 0.431 & 0.7 / -0.2 & 0.07 / 0.00 & 1.08 / 1.00 & yes / yes\\
CIFAR MO AGAD & 4 (10$\times$128) & 600 & 92.1 & 0.027 & 1.8 / 0.3 & 0.18 / 0.00 & 1.20 / 1.02 & no / yes\\
CIFAR MO AGAD & 3 (128$\times$1024) & 2000 & 15.0 & 0.700 & 1.3 / -0.0 & 0.03 / 0.00 & 1.07 / 1.00 & yes / yes\\
CIFAR MO AGAD & 4 (10$\times$128) & 2000 & 96.5 & 0.004 & 32.7 / 0.6 & 0.37 / 0.00 & 1.22 / 1.00 & no / yes\\
LM Prot. c-TTv2 & 0 (128$\times$256) & 750 & 7.6 & 0.717 & 0.3 / 0.1 & 0.01 / 0.00 & 1.00 / 1.00 & yes / yes\\
LM Prot. c-TTv2 & 9 (256$\times$128) & 750 & 0.0 & 0.972 & -1.4 / -0.5 & 0.00 / 0.00 & 0.71 / 1.01 & yes / yes\\
LM Prot. c-TTv2 & 0 (128$\times$256) & 2250 & 32.5 & 0.547 & -0.1 / 0.0 & 0.04 / 0.00 & 1.02 / 1.00 & yes / yes\\
LM Prot. c-TTv2 & 9 (256$\times$128) & 2250 & 0.0 & 0.882 & 0.2 / 0.3 & 0.00 / 0.00 & 0.73 / 1.01 & yes / yes\\
MNIST Li c-TTv2 & 0 (256$\times$784) & 1000 & 21.6 & 0.759 & 0.1 / -0.0 & 0.03 / 0.00 & 1.05 / 1.00 & yes / yes\\
MNIST Li c-TTv2 & 2 (10$\times$128) & 1000 & 3.7 & 0.990 & -0.0 / -0.0 & 0.00 / 0.00 & 1.26 / 0.98 & yes / yes\\
MNIST Li c-TTv2 & 0 (256$\times$784) & 3500 & 25.7 & 0.696 & -0.3 / -0.0 & 0.01 / 0.00 & 1.09 / 1.00 & yes / yes\\
MNIST Li c-TTv2 & 2 (10$\times$128) & 3500 & 5.0 & 0.922 & -1.1 / 0.1 & 0.00 / 0.00 & 1.55 / 1.00 & yes / yes\\
\bottomrule\end{tabular}\end{table*}

